\chapter{Self-contained C1 source definitions and contact verification}
\label{component:p0acdb66bba2f}
\noindent\textbf{Component }\path{p0acdb66bba2f}\par
{\small\noindent Source: \path{CORE/prior/prior_core/c1/source_definitions.tex}\par}
{\small\noindent SHA256: \path{0acdb66bba2fea9b217d8089bcf1df5a1fb50380b01a1a693bbb1b9b1295f241}\par}
{\small\noindent Inventory use: SD1--SD11; full quadratic normal block | SD1-SD11; especially SD2-SD5 | C1\_\allowbreak{}BARE\_\allowbreak{}OTHER\_\allowbreak{}TERMINALS; precise subordinate uses retained below\par}
\medskip
\section{Self-contained C1 source definitions and contact verification}
\label{p0acdb66bba2f:app:c1-source-definitions}

The letter $d$ in the contact--conductor proof denotes the new parameter
$\delta$, not the Bell coefficient. Here the latter is denoted $d_{\rm B}$,
and the original event coefficient $e_5$ is denoted $D_{\rm e}$.
The testing factor is always $D=1+u^2$. The conductor monomial
$C=\delta b^2uw$ is distinct from the clearing multiplier
$\mathcal C=K^6Q^4T_n$.

\paragraph{Domain and seven native coquotients.}
The coefficient field and testing field are independent real closed
fields. On this C1 source branch, $k,f,g>0$ are coefficient
infinitesimals, $e=kf/g$ is positive unbounded, and
$\Delta=e^{-1}$. The shapes $a,t$ are positive bounded elements;
\[
 R=1+k^2(1-af^2),\qquad T=1+k^2-k^2g^2t,
 \qquad M=R-t>0,\quad J=T-a>0.
\]
All original event coefficients are assumed nonnegative. In particular,
the formulas below retain the requirements $D_{\rm e},\beta\ge0$.
No relation compares a testing rate with a coefficient rate.
Put
\[
 \kappa=k^2,\quad \eta=f^2,\quad \zeta=g^2,
 \quad V_c=\kappa\eta\zeta,\quad
 L_c=1-a\eta,\quad H_c=1+(1-a)\eta,\quad Q_c=1-\kappa\zeta+aV_c.
\]
Define
\begin{align*}
P_4={}&1+\kappa+\eta\kappa(a-\kappa-1)-\zeta\kappa t
 +V_c[-\zeta\kappa t^2+2\kappa t+a(1-t)+t],\\
P_\alpha={}&1+\kappa-a\eta\kappa
 +\zeta[(1+\kappa)-(2+\kappa)t]
 +V_c[-\zeta t^2+a(t-1)+t],\\
P_d={}&(1+\kappa)^2-a\eta\kappa^2
 +\zeta\kappa[(\kappa+1)-(\kappa+2)t]\\
 &+V_c[-\zeta\kappa t^2+\kappa a(t-1)+\kappa t+\kappa+1].
\end{align*}
The complete seven coquotients are
\begin{align*}
c_2&=2(1+\kappa)(1-t\zeta)L_c(1+V_c)H_cQ_c,\\
c_b&=2(1+\kappa)L_c[1+(1-t)\zeta]H_cQ_c,\\
c_3&=2(1+\kappa)L_cH_cQ_c(1-\kappa\eta+tV_c),\\
c_5&=2(1+\kappa)L_cH_cQ_c(1+V_c),\\
c_4&=2Q_cP_4,\qquad c_\alpha=2L_cH_cP_\alpha,
 \qquad c_d=4L_cQ_cP_d.
\end{align*}
Their residues, in this order, are $2,2,2,2,2,2,4$.
Here $c_b$ is the original $c_\beta$, and the subscript on $c_d$ is a
name rather than the conductor parameter.

\paragraph{Original source and full parent.}
The original event coefficients and multiplier are
\[
\begin{gathered}
\lambda=n=0,\qquad c=2,\qquad
 d_{\rm B}=\frac{2a\eta c_d}{c_3},\qquad
 m=\frac{2ae^2c_2}{Mc_3},\qquad A=\frac{2a\eta^2c_4}{c_3},\\
\alpha=\frac{2a\Delta^2k^6\eta^2c_\alpha}{c_3},\qquad
D_{\rm e}+\frac{d_{\rm B}}2=\frac{2ac_5}{Jc_3},\qquad
\beta+\frac{d_{\rm B}}2=\frac{2ae^2c_b}{tc_3},\qquad
\Lambda=\frac{RTc_3}{2a}.
\end{gathered}                                                    \tag{SD1}
\]
Let $F$ denote the original complete parent, including its signed core.
For a testing state $(W,X,Y,Z)$ define
\[
\begin{gathered}
 D=1+u^2,\quad D_v=1+v^2,\quad B=X+uZ,\quad Q_A=W+uY,\\
 r=uW-Y,\quad r_2=vQ_A-B,\quad r_5=r-v(Z-uX).
\end{gathered}
\]
With $H_2=4Tc_2$, $H_b=4Tc_b$, $H_3=4Rc_3$, $H_5=4Rc_5$,
the exact entire parent is
\begin{align*}
\Phi:=\Delta^2\Lambda F={}&H_2\frac RM r_2^2
 +H_b\frac RtD_v B^2+\Delta^2H_3\frac TaD_vr^2
 +\Delta^2H_5\frac TJr_5^2\\
 &+4RT\Delta^2\eta^2c_4DD_vY^2
 +4RT\Delta^4k^6\eta^2c_\alpha D(vY-Z)^2\\
 &+2RT\Delta^2\eta c_dD\{(vY-Z)^2-D_vZ^2\}.
                                                               \tag{SD2}
\end{align*}
Thus every original source event is accounted for.

\paragraph{Retained packets and state coordinates.}
In the remaining chart $Z=1$ with $u,v,W,X,Y$ infinitesimal, put
\[
 p=t/R,\quad h=a/T,\quad q_A=v(W+uY),\quad q_B=v(1-uX),
 \quad N=B-pq_A,\quad M_B=r-hq_B.
\]
Retain the four nonnegative packets
\[
 \Pi_1=H_2\frac tM N^2,\quad
 \Pi_2=H_b\frac MtD_vN^2,\quad
 \Pi_3=\Delta^2H_3\frac JaD_vM_B^2,\quad
 \Pi_4=\Delta^2H_5\frac aJM_B^2.                         \tag{SD3}
\]
Their subtraction is the exact identity
\begin{align*}
G=\Phi-\sum_j\Pi_j={}&H_2\{N^2-2q_AN+(1-p)q_A^2\}
 +H_bD_v\{N^2+2q_AN+pq_A^2\}\\
 &+\Delta^2H_3D_v\{M_B^2+2q_BM_B+hq_B^2\}
 +\Delta^2H_5\{M_B^2-2q_BM_B+(1-h)q_B^2\}
 +\mathcal W,\\
\mathcal W={}&4RT\Delta^2\eta^2c_4DD_vY^2
 +4RT\Delta^4k^6\eta^2c_\alpha D(vY-1)^2\\
 &+2RT\Delta^2\eta c_dD\{(vY-1)^2-D_v\}.
                                                               \tag{SD4}
\end{align*}
The testing coordinate $w=(W+uY)/D$ has the full inverse
\[
\begin{gathered}
 X=N+pDvw-u,\qquad q_B=v(D-uN-puDvw),\\
 W=w+\frac{u(M_B+hq_B)}D,\qquad
 Y=uw-\frac{M_B+hq_B}D.
\end{gathered}                                \tag{SD5}
\]
Its only testing denominator is the positive unit $D$.
On the negative-core branch, sign gauges give $u,v>0$, and $w>0$:
after eliminating the two positive core squares the remaining scalar is
\[
 \frac{uvDD_v}{\mathcal D}\{Duv w^2-2DD_vw+uv\},\qquad
 \mathcal D=u^2v^4+u^2v^2+u^2+v^4+2v^2+1>0,
\]
which is nonnegative if $w\le0$. The eliminated matrix is
$\left(\begin{smallmatrix}1+v^2+u^2v^2&uv\\uv&(D+v^2)/D\end{smallmatrix}\right)\succ0$.
The nonnegative-core branch and literal axes already give direct SOS.
Thus the radicals in the conductor proof separate between the two fields.

\paragraph{Coefficient substitution and complete first normals.}
Put
\[
 K=1+k^2,\quad \delta=\Delta R/K,\quad b^2=\eta K/R,
 \quad Q=K+ak^2b^2,\quad T_n=K^2-t\delta^2k^4Qb^2.
\]
Then $\delta,k,b$ are positive infinitesimals, and
\[
 \eta=Kb^2/Q,\quad \Delta=\delta Q/K,\quad
 \zeta=\delta^2k^2Qb^2/K,\quad R=K^2/Q,\quad T=T_n/K.     \tag{SD6}
\]
All displayed denominators are positive pure-coefficient units.
Set $C_i=K^2Q^3c_i$ after this substitution. Explicitly, a native monomial
$c\,k^jf^{2r}g^{2s}a^\alpha t^\tau$ becomes
\[
 c\,\delta^{2s}k^{j+2s}b^{2(r+s)}a^\alpha t^\tau
       K^{2+r-s}Q^{3-r+s}.                            \tag{SD7}
\]
The native coquotients have $-2\le r-s\le3$, so these are polynomials
with no negative exponents or shape inverses.
Define $V_n,Y_n,B_0,A_c$ as in the conductor proof and
$\mathcal C=K^6Q^4T_n$. The factored scalar printed there is exactly
$F_0=(\mathcal CG)|_{N=M_B=0}$. For completeness put
\[
 S_n=8\delta^2K^2b^4C_4D_vY_n
 +8\delta^4Q^2k^6b^4C_\alpha vB_0
 +4\delta^2KQb^2C_dvB_0.
\]
The complete coefficients on the hard normals $N,\delta M_B$ are
\begin{align*}
 L_1={}&8K^3QT_n^2(C_bD_v-C_2)Dvw
       -8\delta^2K^2Q^2uD v^2V_nA_c+K^3auvS_n,\\
 L_2={}&8\delta K^2Q^2T_nD vV_n(C_3D_v-C_5)-K^2T_nS_n/\delta.
                                                               \tag{SD8}
\end{align*}
The quotient in the last line is polynomial because $S_n$ contains
$\delta^2$. These formulas follow by substituting (SD5) into (SD4),
channel by channel; they include all normal terms of degree one.

\paragraph{Short native proof of the exact contact identities.}
Write $V=1+V_c$ and $K=1+\kappa$ in this paragraph. The factored
coquotients above give immediately
\[
 c_b-c_2=\zeta c_3,\qquad c_3-c_5=-\kappa\eta c_b,
 \qquad c_2=(1-t\zeta)c_5.                            \tag{SD9}
\]
Hence $Vc_b-c_2=\zeta c_5$ and
$Vc_3-c_5=-\kappa\eta c_2$.
The remaining three balances are the following short polynomial
identities; substituting the displayed $P_4,P_\alpha,P_d$ proves them
by multiplication and collecting terms:
\begin{align*}
\mathfrak Y:={}&2\kappa c_2+2VL_cc_4
                 -2\kappa^2\zeta Q_cc_\alpha-Q_cc_d=0,\\
\mathfrak S:={}&\kappa L_cc_2+c_5+VL_c^2c_4+\kappa Q_c^2c_\alpha
                 +\tfrac12(\kappa\zeta L_c^2-2L_c-\eta)c_d=0,\\
\mathfrak V:={}&\kappa Rc_2+t\kappa V_cc_b+aV_cc_3+Tc_5
 +KV_cL_c^2c_4-\kappa^2\zeta KQ_cL_cc_\alpha
 -\tfrac K2(Q_cL_c+\eta)c_d=0.                        \tag{SD10}
\end{align*}
For example, after the common factor $4L_cQ_c$ is removed,
the first balance is just
\[
 \kappa K(1-t\zeta)VH_c+VP_4
       -\kappa^2\zeta H_cP_\alpha-Q_cP_d=0.
\]
The two other balances are equally direct identities in the same three
short polynomials; no expanded scalar or ideal calculation enters.

The exact contact, in the original coordinates of (SD2), is
\[
 u=k,\quad v=\Delta k^2\eta,\qquad
 (W,X,Y,Z)=(\Delta k,-k(1-t\zeta),\Delta k^2L_c,1).      \tag{SD11}
\]
At this point
\[
 r_2=Mk\zeta,\quad B=tk\zeta,\quad r=av,\quad
 r_5=-Jv,\quad vY-1=-Q_c,\quad D=K,\quad D_v=V.
\]
Set $\Psi=\Phi/(4RT)$ and let $E_s=\tfrac12\partial_s\Psi$,
with all other original testing coordinates held fixed.
Direct differentiation of the seven terms in (SD2), followed by (SD9),
gives at (SD11)
\begin{align*}
 \Psi&=K\zeta^2\mathfrak S,\\
 E_W&=kv\{\zeta c_2+\Delta^2(Vc_3-c_5)\}=0,\\
 E_X&=k\zeta\{Vc_b-c_2-\zeta c_5\}=0,\\
 E_Y&=\tfrac12K\Delta^3\kappa\eta^2\mathfrak Y,\\
 E_u&=k\zeta^2\{\mathfrak S-c_2+(1-t\zeta)c_5\},\\
 E_v&=v\Delta^2\mathfrak V.
\end{align*}
They all vanish by (SD9)--(SD10). Homogeneity in $(W,X,Y,Z)$ also gives
$E_Z=0$. This proves contact and full first-order stationarity directly
in the native source.

In (SD5)--(SD6), this point is exactly
\[
 (u,w,v,N,M_B)=(k,\delta k,\delta k^2b^2,0,0).
\]
All four packets (SD3) and their first derivatives vanish there.
Changing coordinates and multiplying by $\mathcal C$ therefore gives
$(F_0)_c=(\partial_uF_0)_c=(\partial_wF_0)_c
=(\partial_vF_0)_c=(L_1)_c=(L_2)_c=0$.
The leading polynomial $P$ is a sum of circuit squares vanishing at
this contact, so $P\in J^2$. This proves every identity (CC7) without
using the expanded-source checker as a premise.

\paragraph{Positive correction weight without a large expansion.}
Use the weights $(2,2,6,0,0,2,4,18)$ on
$(\delta,k,b,a,t,u,w,v)$. Write $O_{\rm wt}(s)$ for a polynomial all
of whose monomials have weight at least $s$. Native $\eta$ and $\zeta$
have weights at least $12$ and $20$. Setting them to zero in the seven
short coquotients, and then applying (SD7), gives
\[
 C_i=2K^6+O_{\rm wt}(12)\quad(i=2,b,3,5,4,\alpha),\qquad
 C_d=4K^7+O_{\rm wt}(12).
\]
Also
\[
 Q=K+O_{\rm wt}(16),\quad T_n=K^2+O_{\rm wt}(24),\quad
 V_n=K^2+O_{\rm wt}(24),\quad D_v=1+O_{\rm wt}(36),
\]
\[
 Y_n=O_{\rm wt}(6),\quad B_0=-K^3+O_{\rm wt}(24),\quad
 A_c=2K^8+O_{\rm wt}(12),\quad
 \mathcal C=K^{12}+O_{\rm wt}(16).
\]
In the five factored scalar terms, discard only coefficients carrying
weight at least $12$. The surviving expression is exactly
\[
 8K^{14}D\{D(A^2+B^2)+K(C^2+D_0^2)-2K^2BC\}
       =8K^{14}DP.
\]
Here $A^2,D_0^2$ have weight $44$, while $B^2,C^2,BC$ have weight $40$.
Consequently $E_0=O_{\rm wt}(52)$, sufficient for every scalar lift
coefficient to have weight at least $8$.
For the normals use (SD9) in its polynomial forms
\[
 K(C_b-C_2)=\delta^2k^2Qb^2C_3,\qquad
 Q(C_3-C_5)=-k^2Kb^2C_b.
\]
Since $K,Q$ have constant term one, these show weights at least $20,16$
for the two differences, respectively, without putting their inverses
in any lifting coefficient. The displayed $S_B,S_C$ in the conductor
proof have weight at least $12$. Its factored normal formulas now give
$L_1=O_{\rm wt}(42)$ and $L_2=O_{\rm wt}(32)$.
The monic lift therefore has positive coefficient weights at least
$8,20,10$. The stronger exact checker bounds $12,20,12$ are optional.

\paragraph{Full hard block and restoration.}
There are no state-normal terms above degree two. Put
\[
 C_B=H_3D_vh+H_5(1-h),\qquad
 B_Y=4RT\Delta^2\eta^2c_4DD_v
 +4RT\Delta^4k^6\eta^2c_\alpha Dv^2
 +2RT\Delta^2\eta c_dDv^2.
\]
Before multiplication by $\mathcal C$, the exact quadratic block on
$(N,\delta M_B)$ has entries
\begin{align*}
 A_{11}&=H_2+H_bD_v+\Delta^2C_Bu^2v^2+B_Yh^2u^2v^2/D^2,\\
 A_{12}&=-\frac{\Delta^2}{\delta}(H_3D_v-H_5)uv
                  -\frac{B_Y}{\delta D^2}huv,\\
 A_{22}&=\frac{\Delta^2}{\delta^2}(H_3D_v+H_5)
                  +\frac{B_Y}{\delta^2D^2}.
\end{align*}
The displayed ratios cancel coefficient monomials exactly:
$\Delta/\delta=Q/K$ is a positive unit and
$B_Y/\delta^2$ is bounded infinitesimal. Thus this full block has
residue $16I_2$, also after multiplication by $\mathcal C$.
The conductor proof and the positive circuit frame give the bounded
strict-residue Gram for $\mathcal CG$. Restore the four full positive
packets $\mathcal C\Pi_j$ from (SD3), then undo
$\mathcal C\Delta^2\Lambda$. This recovers the entire original parent;
retained gap ratios need not be bounded and are never placed in that
bounded Gram. The complementary state and outer-angle charts require
the already proved C1 chart arguments. This appendix makes no claim of
an exhaustive global source cover.
